\documentclass[border=5mm]{standalone}
% ==============================================================================
% SETUP.TEX - CONFIGURACIÓN GLOBAL DEL PROYECTO
% ==============================================================================
% Este archivo centraliza todos los paquetes y estilos.
% Debe incluirse en el preámbulo del libro principal y de los archivos standalone.
% \PassOptionsToPackage{gray}{xcolor}
% --- 1. CODIFICACIÓN E IDIOMA ---
% fix-cm habilita las fuentes Computer Modern en tamaños arbitrarios (por debajo
% de 5 pt, entre otros). Debe cargarse antes que fontenc.
\usepackage{fix-cm}
\usepackage[utf8]{inputenc}
\usepackage[T1]{fontenc}
\usepackage[spanish, es-tabla]{babel}  
\usepackage{csquotes} % Recomendado para babel y citas

\usepackage{import}
% mode=image: Usa el PDF pre-compilado, nunca intenta recompilar.
% Debes compilar cada figura manualmente antes de compilar el libro.
\usepackage[mode=image, subpreambles=false]{standalone}

% --- 2. MATEMÁTICAS Y CIENCIAS ---
\usepackage{amsmath, amssymb, amsfonts, mathtools}
\usepackage{enumitem}

% LaTeX trae \DeclareMathSizes{5}{5}{5}{5}, así que a 5 pt (\tiny) los subíndices
% salen del mismo tamaño que la base: en las etiquetas de figuras el M de
% $\omega_M$ queda desproporcionado. Se restablece la proporción habitual.
\DeclareMathSizes{5}{5}{3.7}{3}

% --- 3. DISEÑO Y GEOMETRÍA --- 
% Unificamos los márgenes para todo el documento
\usepackage[left=2.5cm,right=2.5cm,top=3cm,bottom=3cm]{geometry}
\makeatletter
\ifstandalone % Evita que geometry dañe el recorte de las figuras bajándolas 3cm
  \@ifclasswith{standalone}{crop=false}{
    % Es un capítulo/sección: Dejamos que geometry actúe.
  }{
    % Es una figura: Neutralizamos geometry para que no las recorte mal.
    \geometry{pass}
  }
\fi
\makeatother
\usepackage{parskip} % Espaciado entre párrafos sin sangría
\usepackage{float}   % Para usar [H] en figura

% Ajustes de flotantes para evitar espacios verticales exagerados
\renewcommand{\topfraction}{0.95}      % Permite que las figuras ocupen hasta el 95% superior
\renewcommand{\bottomfraction}{0.95}   % Permite que las figuras ocupen hasta el 95% inferior
\renewcommand{\textfraction}{0.05}     % Permite hojas con tan solo un 5% de texto
\renewcommand{\floatpagefraction}{0.8} % Requiere que una página de solo figuras esté 80% llena
\setcounter{topnumber}{4}
\setcounter{bottomnumber}{4}
\setcounter{totalnumber}{10}

% --- 4. GRÁFICOS Y TABLAS ---
\usepackage{graphicx}
\usepackage{booktabs} % Tablas profesionales
\usepackage{caption}  % Control de captions y \captionof
\usepackage{tikz}
\usetikzlibrary{arrows.meta, calc, angles, quotes, babel, backgrounds}
\usepackage{pgfplots}
\usepgfplotslibrary{groupplots}
\usepgfplotslibrary{external}
\usepgfplotslibrary{patchplots}
\pgfplotsset{compat=1.18}
\usepackage[american, straightvoltages]{circuitikz}

% --- 5. COLORES Y CAJAS ---

\usepackage[table]{xcolor}
\usepackage{tcolorbox}

% Definición de colores corporativos/estilo
\definecolor{utnblue}{RGB}{0, 51, 102}
\definecolor{codegreen}{rgb}{0,0.6,0}
\definecolor{codegray}{rgb}{0.5,0.5,0.5}
\definecolor{codepurple}{rgb}{0.58,0,0.82}
\definecolor{backcolour}{rgb}{0.96,0.96,0.96}

% --- 6. ENTORNOS PERSONALIZADOS (CAJAS) ---

% Caja "Resultado" (Azul Corporativo) — Definiciones, fórmulas clave, resultados importantes.
% Uso: \begin{resultado}[Fórmula de Euler] ... \end{resultado}
\newtcolorbox{resultado}[1][]{
    colback=utnblue!5!white,
    colframe=utnblue,
    title=\textbf{#1},
    fonttitle=\bfseries,
    sharp corners,
    boxrule=0.5mm
}

% Caja "Nota" (Gris) — Advertencias, aplicaciones, contexto secundario.
% Uso: \begin{nota}[Cuidado con la calculadora] ... \end{nota}
% Uso: \begin{nota} ... \end{nota}  → título por defecto "Nota"
\newtcolorbox{nota}[1][Nota]{
    colback=codegray!5!white,
    colframe=codegray,
    title=\textbf{#1},
    fonttitle=\bfseries,
    sharp corners,
    boxrule=0.5mm
}

% Caja inline para resaltar un valor y enlazarlo con su otra aparición en una tabla
% o en una cuenta: mismo color, mismo valor.
% Uso: \marcavalor{$4$}  o  \marcavalor[codegreen]{$4$}
\newtcbox{\marcavalor}[1][utnblue]{
    on line,
    boxsep=1pt, left=2pt, right=2pt, top=1pt, bottom=1pt,
    colback=#1!10!white,
    colframe=#1,
    boxrule=0.4pt,
    arc=2pt
}

% --- 7. CÓDIGO FUENTE (LISTINGS) ---
\usepackage{listings}

\lstdefinestyle{miestilo}{
    backgroundcolor=\color{backcolour},   
    commentstyle=\color{codegreen},
    keywordstyle=\color{magenta},
    numberstyle=\tiny\color{codegray},
    stringstyle=\color{codepurple},
    basicstyle=\ttfamily\footnotesize,
    breakatwhitespace=false,         
    breaklines=true,                 
    captionpos=b,                    
    keepspaces=true,                 
    numbers=left,                    
    numbersep=5pt,                  
    showspaces=false,                
    showstringspaces=false,
    showtabs=false,                  
    tabsize=2,
    frame=single,                    
    rulecolor=\color{codegray}
}

\lstset{style=miestilo}
% Soporte para tildes y ñ en bloques de código
\lstset{literate={ñ}{{\~n}}1 {á}{{\'a}}1 {é}{{\'e}}1 {í}{{\'i}}1 {ó}{{\'o}}1 {ú}{{\'u}}1}

% Operadores matemáticos personalizados

% Vector en sentido discreto (lista finita de coordenadas): negrita + flecha.
% Uso: \vect{v}, \vect{e}_k, \vect{0}.
\newcommand{\vect}[1]{\vec{\mathbf{#1}}}

% Fecha de versión y comando para capítulos standalone
\providecommand{\verdate}{\the\year.\ifnum\month<10 0\fi\the\month.\ifnum\day<10 0\fi\the\day}
\newcommand{\chapterversion}{%
    \vspace{-1.5cm}\begin{flushright}\small\textit{\color{codegray}Versión~\verdate}\end{flushright}\vspace{0.5cm}%
}

% --- 8. HIPERVÍNCULOS (DEBE SER EL ÚLTIMO PAQUETE) ---
\usepackage[colorlinks=true, linkcolor=utnblue, urlcolor=utnblue, citecolor=utnblue]{hyperref}

% --- 9. REFERENCIAS CRUZADAS ENTRE CAPÍTULOS STANDALONE ---
% Cuando se compila un capítulo standalone, xr-hyper lee los .aux de los
% otros capítulos (si existen) para resolver \ref{} a labels externos.
% En la compilación del libro completo este bloque no se activa.
\ifstandalone
  \usepackage{xr-hyper}
  \makeatletter
  \newcommand{\xrchap}[1]{%
    \edef\xrc@self{\detokenize{Capitulo#1}}%
    \edef\xrc@job{\jobname}%
    \ifx\xrc@self\xrc@job\else
      \IfFileExists{../#1capitulo/Capitulo#1.aux}%
        {\externaldocument{../#1capitulo/Capitulo#1}}{}%
    \fi
  }
  \makeatother
  \xrchap{01}\xrchap{02}\xrchap{03}\xrchap{04}
  \xrchap{05}\xrchap{06}\xrchap{07}\xrchap{08}
  \xrchap{09}\xrchap{10}\xrchap{11}\xrchap{12}
  \xrchap{13}
\fi
% \PassOptionsToPackage{gray}{xcolor}

% fig_zoh
% (a) tiempo: escalera (const plot) sobre muestras vs senal ideal (punteada).
% (b) frecuencia: copias en 0,+-ws,+-2ws (ws=3) con la envolvente seno cardinal
%     |sinc(w*pi/3)| (Delta t = 2pi/ws) por encima, ceros en multiplos de ws.
% (c) el producto: cada copia multiplicada por la envolvente, con el filtro de
%     reconstruccion marcado (corte en ws/2 = 1.5 y realce 1/|sinc| dentro de la banda).
% (d) el espectro final: sobrevive solo la banda base, sin imagenes y con la caida
%     de la envolvente ya compensada por el filtro.
% El +1e-6 de los argumentos evita la division por cero de sinc en w=0 y en w=+-ws.

\begin{document}
\begin{tikzpicture}[>={Latex[length=3pt]}, font=\scriptsize]

% ===================== (a) tiempo: escalera =====================
\begin{axis}[
    name=pa,
    width=7.6cm, height=4.6cm,
    axis lines=center, axis line style={->, thick},
    xmin=-0.4, xmax=6.3, ymin=-0.5, ymax=1.2,
    xtick={0,1,2,3,4,5}, xticklabels={}, ytick=\empty,
    xlabel={$t$}, xlabel style={anchor=west, at={(axis description cs:1.0,0.42)}},
    title={(a)}, title style={font=\scriptsize, yshift=-3pt},
    enlarge x limits=false, clip mode=individual, axis on top,
]
    % senal ideal (punteada)
    \addplot[gray!70, dotted, thick, smooth] coordinates {
        (0,0.30) (1,0.70) (2,1.00) (3,0.75) (4,0.30) (5,-0.20)
    };
    % escalera ZOH (mantiene v_n sobre [n,n+1))
    \addplot[const plot, utnblue, thick] coordinates {
        (0,0.30) (1,0.70) (2,1.00) (3,0.75) (4,0.30) (5,-0.20) (6,-0.20)
    };
    % muestras
    \addplot[only marks, mark=*, mark size=1.5pt,
             mark options={fill=black, draw=black}]
        coordinates {(0,0.30) (1,0.70) (2,1.00) (3,0.75) (4,0.30) (5,-0.20)};
\end{axis}

% ===================== (b) frecuencia: envolvente seno cardinal =====================
\begin{axis}[
    name=pb,
    at={($(pa.south east)+(1.5cm,0)$)}, anchor=south west,
    width=7.6cm, height=4.6cm,
    axis lines=center, axis line style={->, thick},
    xmin=-7.2, xmax=7.4, ymin=-0.12, ymax=1.05,
    xtick={-6,-3,3,6}, xticklabels={$-2\omega_s$,$-\omega_s$,$\omega_s$,$2\omega_s$},
    ytick=\empty,
    every tick label/.style={font=\tiny},
    xlabel={$\omega$}, xlabel style={anchor=west, at={(axis description cs:1.0,0.18)}},
    title={(b)}, title style={font=\scriptsize, yshift=-3pt},
    enlarge x limits=false, clip mode=individual, axis on top,
]
    % copias (triangulos) en 0,+-3,+-6
    \addplot[gray!55, thick, fill=gray!10] coordinates {(-1,0)(0,0.85)(1,0)};
    \addplot[gray!55, thick, fill=gray!10] coordinates {(2,0)(3,0.85)(4,0)};
    \addplot[gray!55, thick, fill=gray!10] coordinates {(-4,0)(-3,0.85)(-2,0)};
    \addplot[gray!55, thick, fill=gray!10] coordinates {(5,0)(6,0.85)(7,0)};
    \addplot[gray!55, thick, fill=gray!10] coordinates {(-7,0)(-6,0.85)(-5,0)};
    % envolvente seno cardinal |Delta t sinc(w Delta t/2)|, Delta t=2pi/3
    \addplot[red!80!black, thick, samples=400, domain=-6.95:7.02]
        {abs(sin(deg(x*pi/3))/(x*pi/3))*0.85};
    \node[red!80!black, anchor=west, font=\tiny] at (axis cs:1.35,0.86) {envolvente};
\end{axis}

% ===================== (c) el producto y el filtro de reconstruccion =====================
\begin{axis}[
    name=pc,
    at={($(pa.south west)+(0,-1.0cm)$)}, anchor=north west,
    width=7.6cm, height=4.6cm,
    axis lines=center, axis line style={->, thick},
    xmin=-7.2, xmax=7.4, ymin=-0.12, ymax=1.05,
    xtick={-6,-3,-1.5,1.5,3,6},
    xticklabels={$-2\omega_s$,$-\omega_s$,,$\tfrac{\omega_s}{2}$,$\omega_s$,$2\omega_s$},
    ytick=\empty,
    every tick label/.style={font=\tiny},
    xlabel={$\omega$}, xlabel style={anchor=west, at={(axis description cs:1.0,0.18)}},
    title={(c)}, title style={font=\scriptsize, yshift=-3pt},
    enlarge x limits=false, clip mode=individual, axis on top,
]
    % copias multiplicadas por la envolvente
    \addplot[gray!55, thick, fill=gray!10, samples=120, domain=-1:1]
        {0.85*(1-abs(x))*abs(sin(deg((x+1e-6)*pi/3))/((x+1e-6)*pi/3))} \closedcycle;
    \addplot[gray!55, thick, fill=gray!10, samples=120, domain=2:4]
        {0.85*(1-abs(x-3))*abs(sin(deg((x+1e-6)*pi/3))/((x+1e-6)*pi/3))} \closedcycle;
    \addplot[gray!55, thick, fill=gray!10, samples=120, domain=-4:-2]
        {0.85*(1-abs(x+3))*abs(sin(deg((x+1e-6)*pi/3))/((x+1e-6)*pi/3))} \closedcycle;
    \addplot[gray!55, thick, fill=gray!10, samples=120, domain=5:7]
        {0.85*(1-abs(x-6))*abs(sin(deg((x+1e-6)*pi/3))/((x+1e-6)*pi/3))} \closedcycle;
    \addplot[gray!55, thick, fill=gray!10, samples=120, domain=-7:-5]
        {0.85*(1-abs(x+6))*abs(sin(deg((x+1e-6)*pi/3))/((x+1e-6)*pi/3))} \closedcycle;
    % filtro de reconstruccion: recuadro con corte en ws/2, como el pasabajos ideal
    % de fig_reconstruccion_lp
    \addplot[red!80!black, thick, dotted]
        coordinates {(-1.5,0) (-1.5,0.95) (1.5,0.95) (1.5,0)};
    \node[red!80!black, anchor=west, font=\tiny] at (axis cs:1.62,0.90) {filtro};
\end{axis}

% ===================== (d) el espectro final =====================
\begin{axis}[
    name=pd,
    at={($(pb.south west)+(0,-1.0cm)$)}, anchor=north west,
    width=7.6cm, height=4.6cm,
    axis lines=center, axis line style={->, thick},
    xmin=-7.2, xmax=7.4, ymin=-0.12, ymax=1.05,
    xtick={-6,-3,3,6}, xticklabels={$-2\omega_s$,$-\omega_s$,$\omega_s$,$2\omega_s$},
    ytick=\empty,
    every tick label/.style={font=\tiny},
    xlabel={$\omega$}, xlabel style={anchor=west, at={(axis description cs:1.0,0.18)}},
    title={(d)}, title style={font=\scriptsize, yshift=-3pt},
    enlarge x limits=false, clip mode=individual, axis on top,
]
    % banda base recuperada, sin imagenes y con la caida de la envolvente compensada
    \addplot[utnblue, thick, fill=utnblue!14] coordinates {(-1,0)(0,0.85)(1,0)};
\end{axis}
\end{tikzpicture}
\end{document}
